% ============================================================================
\chapter{Supervised Learning from an Engineer's Perspective}
\label{ch:learning}
% ============================================================================

This chapter introduces the concepts of supervised learning---the
framework used to train neural networks from labelled data. If you
have a machine learning background, you can skim this chapter; it
establishes the vocabulary and intuition needed for the training
methods in Parts~II--IV.


% ----------------------------------------------------------------------------
\section{Classification as Function Approximation}
\label{sec:classification}
% ----------------------------------------------------------------------------

A \concept{classifier} is a function that maps an input (an image,
a sensor reading, a document) to one of several predefined categories.
For MNIST:
\begin{equation}
  f : \{0, 1, \ldots, 255\}^{784} \to \{0, 1, 2, \ldots, 9\}
  \label{eq:classifier}
\end{equation}

We do not know the ``true'' function $f^*$ that perfectly classifies
every possible handwritten digit. Instead, we have a dataset of
examples $({\bvec{x}_i, y_i})$ and we want to find a function
$\hat{f}$ that agrees with these examples and \emph{generalises} to
unseen inputs.

\begin{analogy}
Learning a classifier from data is like learning to identify birds from
a field guide. You study examples (``this is a robin, this is a
sparrow''), and eventually you can identify birds you have never seen
before---not because you memorised every bird, but because you learned
the distinguishing features (beak shape, colouring, size).
\end{analogy}


% ----------------------------------------------------------------------------
\section{Training Data, Test Data, and Overfitting}
\label{sec:overfitting}
% ----------------------------------------------------------------------------

The 70,000 MNIST~\cite{lecun1998mnist} images are split into:
\begin{itemize}
  \item \textbf{Training set} (60,000 images): used to adjust the
    model's parameters.
  \item \textbf{Test set} (10,000 images): used only to evaluate
    accuracy. The model never sees these during training.
\end{itemize}

\concept{Overfitting} occurs when a model memorises the training
examples instead of learning general patterns. An overfitting model
achieves high training accuracy but poor test accuracy.

\begin{keyinsight}
The test set is the only honest measure of a model's quality. Any
accuracy number reported in this book refers to the test set unless
explicitly stated otherwise.
\end{keyinsight}


% ----------------------------------------------------------------------------
\section{The Traditional Artificial Neuron}
\label{sec:neuron}
% ----------------------------------------------------------------------------

The building block of conventional neural networks is the
\concept{artificial neuron}. It computes:

\begin{equation}
  y = \sigma\!\left(\sum_{i=1}^{n} w_i \cdot x_i + b\right)
    = \sigma(\bvec{w} \cdot \bvec{x} + b)
  \label{eq:neuron}
\end{equation}

where:
\begin{itemize}
  \item $\bvec{x} = (x_1, \ldots, x_n)$ is the input vector.
  \item $\bvec{w} = (w_1, \ldots, w_n)$ is the \concept{weight}
    vector (learned parameters).
  \item $b$ is the \concept{bias} (a learned offset).
  \item $\sigma$ is the \concept{activation function} (a nonlinear
    function such as ReLU: $\sigma(z) = \max(0, z)$).
\end{itemize}

\begin{figure}[htbp]
\centering
\begin{tikzpicture}[
  input/.style={circle, draw, minimum size=0.6cm, font=\small},
  neuron/.style={circle, draw, minimum size=1cm, font=\small,
    fill=bitblue!10, line width=0.8pt},
  >=Stealth
]
  \foreach \i/\v in {1/x_1, 2/x_2, 3/x_n} {
    \ifnum\i=3
      \node[input] (in\i) at (0, 3-\i*1.2) {$\v$};
    \else
      \node[input] (in\i) at (0, 3-\i*1.2) {$\v$};
    \fi
  }
  \node at (0, -0.4) {$\vdots$};

  \node[neuron] (n) at (3.5, 0.6) {$\Sigma$};
  \node[draw, rounded corners, fill=bitorange!10, minimum width=0.8cm,
    font=\small] (act) at (5.5, 0.6) {$\sigma$};
  \node[right=0.5cm, font=\small] (out) at (act.east) {$y$};

  \foreach \i in {1,2,3} {
    \draw[->, thick] (in\i) -- node[above, font=\scriptsize, sloped] {$w_\i$} (n);
  }
  \draw[->, thick] (n) -- (act);
  \draw[->, thick] (act) -- (out);
\end{tikzpicture}
\caption{A traditional artificial neuron. The weighted sum
  $\sum w_i x_i + b$ is computed first, then passed through a nonlinear
  activation function $\sigma$.}
\label{fig:neuron}
\end{figure}

The weighted sum $\bvec{w} \cdot \bvec{x} + b$ is the
\concept{multiply-accumulate} (MAC) operation that dominates the cost
of conventional networks. For a neuron with 784 inputs (one per MNIST
pixel), this requires 784 floating-point multiplications and 784
floating-point additions per neuron.


% ----------------------------------------------------------------------------
\section{Loss Functions}
\label{sec:loss}
% ----------------------------------------------------------------------------

To train a model, we need a way to measure \emph{how wrong} its
predictions are. This measurement is called the \concept{loss function}
(or cost function).

For 10-class classification (MNIST), the standard loss is
\concept{cross-entropy loss}. The model outputs a probability
distribution over 10 classes, and we penalise it for assigning low
probability to the correct class.

\begin{equation}
  \mathcal{L} = -\log p_{\hat{y}}
  \label{eq:cross-entropy}
\end{equation}

where $p_{\hat{y}}$ is the model's predicted probability for the
correct class $\hat{y}$.

\begin{example}
Suppose the true label is ``3'' and the model outputs:
\[
  \bvec{p} = (0.01, 0.02, 0.05, \mathbf{0.80}, 0.03, 0.01, 0.02, 0.01, 0.03, 0.02)
\]
The loss is $\mathcal{L} = -\log(0.80) \approx 0.22$.

If the model were less confident, outputting $p_3 = 0.10$, the loss
would be $\mathcal{L} = -\log(0.10) \approx 2.30$---much larger,
reflecting the worse prediction.
\end{example}

\begin{engineernote}
Not all paradigms use cross-entropy loss. Tsetlin Machines
(\cref{ch:clauses}) use a completely different training signal based on
reinforcement learning. Logic Gate Networks (\cref{ch:diff-gates}) use
cross-entropy or similar differentiable losses. Understanding the
concept of ``measuring wrongness'' is the key; the specific formula
varies.
\end{engineernote}


% ----------------------------------------------------------------------------
\section{Gradient Descent}
\label{sec:gradient-descent}
% ----------------------------------------------------------------------------

Given a loss function that measures prediction error, \concept{gradient
descent} is the algorithm that adjusts the model's parameters to reduce
that error.

\begin{enumerate}
  \item \textbf{Forward pass:} Run an input through the model, get a
    prediction.
  \item \textbf{Compute loss:} Measure how wrong the prediction is.
  \item \textbf{Compute gradient:} For each parameter $w$, compute
    $\frac{\partial \mathcal{L}}{\partial w}$---how much the loss
    changes if we nudge $w$ slightly.
  \item \textbf{Update:} Move each parameter in the direction that
    reduces the loss:
    \begin{equation}
      w \leftarrow w - \eta \cdot \frac{\partial \mathcal{L}}{\partial w}
      \label{eq:gradient-update}
    \end{equation}
    where $\eta > 0$ is the \concept{learning rate} (step size).
\end{enumerate}

\begin{analogy}
Imagine you are blindfolded on a hillside and want to reach the valley.
You feel the slope under your feet (the gradient) and take a step
downhill. Repeat. You will eventually reach the bottom---maybe not the
absolute lowest point, but a low one. Gradient descent does the same in
the high-dimensional space of model parameters.
\end{analogy}

\begin{figure}[htbp]
\centering
\begin{tikzpicture}[>=Stealth]
  \begin{axis}[
    width=9cm, height=5cm,
    xlabel={Parameter $w$}, ylabel={Loss $\mathcal{L}$},
    xmin=-3, xmax=3, ymin=0, ymax=5,
    axis lines=middle,
    every axis label/.style={font=\small\sffamily},
    ticks=none
  ]
    % Loss curve
    \addplot[thick, bitblue, smooth, domain=-3:3, samples=100]
      {0.5*x^2 + 0.3*sin(deg(2*x)) + 1};

    % Gradient descent steps
    \addplot[only marks, mark=*, mark size=2.5pt, bitorange]
      coordinates {(2, 3.21) (1.2, 1.92) (0.5, 1.41) (0.1, 1.11)};
    \draw[->, thick, bitorange] (2, 3.21) -- (1.2, 1.92);
    \draw[->, thick, bitorange] (1.2, 1.92) -- (0.5, 1.41);
    \draw[->, thick, bitorange] (0.5, 1.41) -- (0.1, 1.11);

    \node[font=\scriptsize\sffamily, bitorange, right] at (2.1, 3.5)
      {start};
  \end{axis}
\end{tikzpicture}
\caption{Gradient descent iteratively moves the parameter $w$ in the
  direction that reduces the loss. Each step follows the negative
  gradient (downhill).}
\label{fig:gradient-descent}
\end{figure}


% ----------------------------------------------------------------------------
\section{Backpropagation}
\label{sec:backprop}
% ----------------------------------------------------------------------------

In a multi-layer network, computing $\frac{\partial \mathcal{L}}
{\partial w}$ for a parameter deep in the network requires the
\concept{chain rule} of calculus.

If the network computes $y = f(g(h(x)))$, then:
\begin{equation}
  \frac{\partial y}{\partial x} =
    \frac{\partial f}{\partial g} \cdot
    \frac{\partial g}{\partial h} \cdot
    \frac{\partial h}{\partial x}
  \label{eq:chain-rule}
\end{equation}

\concept{Backpropagation} applies this chain rule layer by layer, from
the output back to the input, efficiently computing gradients for every
parameter in the network.

\begin{engineernote}
Backpropagation requires that every operation in the forward pass be
\emph{differentiable}---that is, it must have a well-defined
derivative. This is the central challenge for bitwise networks:
discrete operations like AND and OR are \emph{not} differentiable.
Their output jumps from 0 to~1 with no smooth transition, so the
derivative is zero everywhere (useless for learning) or undefined
(at the jump). The straight-through estimator~\cite{bengio2013ste} is one common
workaround for this problem.
\end{engineernote}


% ----------------------------------------------------------------------------
\section{The Differentiability Problem}
\label{sec:diff-problem}
% ----------------------------------------------------------------------------

Consider the AND function: $f(a, b) = a \wedge b$. The function maps
$\{0,1\}^2$ to $\{0,1\}$---it is defined only at four points. Between
those points, there is no continuous surface, so there is no gradient.

\begin{figure}[htbp]
\centering
\begin{tikzpicture}[>=Stealth]
  % Continuous neuron
  \begin{scope}
    \node[font=\sffamily\bfseries] at (2, 3.2) {Continuous (ReLU)};
    \begin{axis}[
      at={(0,0)}, width=4.5cm, height=3.5cm,
      xmin=-2, xmax=2, ymin=-0.5, ymax=2.5,
      axis lines=middle, ticks=none,
      every axis label/.style={font=\tiny\sffamily}
    ]
      \addplot[very thick, bitblue, domain=-2:2, samples=50]
        {max(0, x)};
      \node[font=\tiny\sffamily, bitblue] at (axis cs:1.5, 2)
        {smooth};
      \draw[->, thick, bitgreen] (axis cs:1, 1) -- (axis cs:1.3, 1)
        node[right, font=\tiny\sffamily, bitgreen] {$\nabla$};
    \end{axis}
  \end{scope}

  % Discrete gate
  \begin{scope}[xshift=6cm]
    \node[font=\sffamily\bfseries] at (2, 3.2) {Discrete (AND gate)};
    \begin{axis}[
      at={(0,0)}, width=4.5cm, height=3.5cm,
      xmin=-0.5, xmax=1.5, ymin=-0.5, ymax=1.5,
      axis lines=middle, ticks=none,
      every axis label/.style={font=\tiny\sffamily}
    ]
      \addplot[only marks, mark=*, mark size=3pt, bitorange]
        coordinates {(0,0) (1,0) (0,0) (1,1)};
      \node[font=\tiny\sffamily, bitorange] at (axis cs:0.8, 0.5)
        {no gradient};
    \end{axis}
  \end{scope}
\end{tikzpicture}
\caption{Left: a continuous activation (ReLU) has a gradient that
  tells us which direction to move. Right: a discrete logic gate has
  no gradient---its outputs are isolated points with no smooth
  connection.}
\label{fig:diff-problem}
\end{figure}

This is the fundamental obstacle that every bitwise neural network must
overcome. The three paradigms in this book handle it in three different
ways:


% ----------------------------------------------------------------------------
\section{Three Solutions to the Differentiability Problem}
\label{sec:three-solutions}
% ----------------------------------------------------------------------------

\begin{description}[leftmargin=1em]
  \item[Tsetlin Machines~\cite{granmo2018tsetlin}: avoid gradients entirely.]
    Instead of computing derivatives, Tsetlin Machines use
    \emph{reinforcement learning}---simple reward/penalty signals that
    nudge automata states up or down. No calculus required. The
    learning rule is: ``if the clause helped classify correctly,
    reinforce it; if it caused an error, weaken it.'' This is covered
    in Chapters~\ref{ch:automaton}--\ref{ch:tm-mnist}.

  \item[Logic Gate Networks~\cite{petersen2022difflogic}: use a continuous relaxation.]
    During training, each node does not commit to a single gate type.
    Instead, it maintains a \emph{probability distribution} over all 16
    gate types, and its output is a differentiable weighted average.
    Gradients flow through this weighted average. As training
    progresses, the distribution sharpens until each node converges to
    a single gate. This is covered in Chapters~\ref{ch:diff-gates}
    --\ref{ch:lgn-training}.

  \item[KAN/LUT~\cite{liu2024kan}: train in continuous space, deploy discretely.]
    During training, learnable functions are smooth B-spline curves~\cite{deboor1978splines}
    that support standard backpropagation. After training, these
    curves are \emph{compiled} into discrete lookup tables. The
    discretisation happens only once, at deployment time. This is
    covered in Chapters~\ref{ch:splines}--\ref{ch:lut}.
\end{description}

\begin{keyinsight}
Each paradigm makes a different trade-off:
\begin{itemize}
  \item Tsetlin: no gradients at all (simplest training, slowest
    convergence).
  \item LGN: gradients through a relaxation (standard training, must
    anneal to discrete).
  \item KAN/LUT: full gradients during training, discretisation at
    deployment (highest training cost, smallest deployment cost).
\end{itemize}
The hybrid architecture uses all three, each where it is most
effective.
\end{keyinsight}
